% Article VII, additive recovery, revision 04, 10-09-2026.
% RESULTADO_RECUPERADO: S15, 11c_gravedad_holonomia_rev6.tex:16--126.
% The canonical cocycle c27=(1,18) is preserved as an antecedent of the route
% computation; this fragment proves its composition and representation.
% Finite translational memory and the cosmological amplitude s0 have
% different types. The cocycle does not by itself evaluate a local density.

\section{Phase return and representation of accumulated memory}
\label{vii:sec:retorno-memoria-gravedad}

The chronology of enriched routes distinguishes positional return,
orientational return, and phase return. The observable representation
identifies certain arrival states, while the cumulative record
preserves the sequence traversed. The geometric representation must
transport both data jointly.

Preservation of the complement in
\eqref{g26:eq:archivo-completo-en} and reconstruction of residuals in
\eqref{g26:eq:accion-residuo-en} use this complete temporal record.
Two paths with the same observable phase may have different memory
cocycles. Transport of the radial reader preserves both coordinates:
periodic return acts on phase, while concatenation increments the record
appearing in the following formulas.

\subsection{Inversion cocycle and change of sheet}
\label{vii:subsec:cociclo-memoria}

Let \(F_{\rm obs}\) be the chronological representation of TPK transport.
Its canonical blocks have the hierarchy
\[
\begin{array}{c|l}
27&\text{positional return with reversal of orientation},\\
54&\text{positional and orientational return},\\
108&\text{complete return of the observable phase}.
\end{array}
\]
The record counts reversals of orientation \(g(a)\) and
effective changes of sheet \(s(a)\). For a route,
\begin{equation}
 c(\gamma)=\sum_{a\subset\gamma}(g(a),s(a)),\qquad
 c(\gamma_2\gamma_1)=c(\gamma_2)+c(\gamma_1).
 \label{vii:eq:cociclo-ruta}
\end{equation}
The canonical computation of the block of length \(27\) expresses
\(c_{27}=(1,18)\). Composition preserves those increments in
the four blocks completing the observable period.
On forward-oriented routes the record is cumulative;
its extension to the route groupoid uses \(\mathbb Z^2\) and
\(c(\bar\gamma)=-c(\gamma)\).

\begin{proposition}[Memory of the complete observable circuit]
\label{vii:prop:memoria-vuelta-completa}
Writing \(\mathcal K_{27}=F_{\rm obs}^{27}\), the lift
to the record, for \(x\in\mathcal O_{108}\) and \(z\in\mathbb Z^2\), satisfies
\[
 \widetilde{\mathcal K}_{27}(x,z)
   =(\mathcal K_{27}x,z+(1,18)),\qquad
 \widetilde{\mathcal K}_{27}^{\,4}(x,z)
   =(x,z+(4,72)).
\]
After \(N\) complete circuits, the increment is \(N(4,72)\).
\end{proposition}

\begin{proof}
The observable projection of \(\mathcal K_{27}\) has order four.
The preceding canonical computation fixes \(c_{27}=(1,18)\) constantly
on that observable orbit. Additivity of
\eqref{vii:eq:cociclo-ruta} sums the four increments and gives
\(c_{108}=4(1,18)=(4,72)\).
Repetition applies the same concatenation law and produces
\(Nc_{108}\). Return of the first coordinate preserves the
increment of the second.
\end{proof}

\subsection{Affine representation of the cocycle}
\label{vii:subsec:representacion-afin-cociclo}

The discrete realization uses an element \(R_4\) of order four in
\(\operatorname{Spin}^+(1,3)\), a unit timelike vector \(u\)
fixed by its vector action, and a length unit \(\ell_0>0\).
The scale \(\ell_0\) preserves its provenance in the dimensional section.
We define
\begin{equation}
 \rho_{\rm C}^{\rm disc}(\gamma)=
 \bigl(R_4^{c_g(\gamma)},\,\ell_0c_s(\gamma)u\bigr)
 \in\operatorname{Spin}^+(1,3)\ltimes\mathbb R^{1,3}.
 \label{vii:eq:representacion-discreta-memoria}
\end{equation}
The power acts on the vector through the associated vector
representation of the spin group.

\begin{theorem}[Composition and translational amplitude of return]
\label{vii:thm:traslacion-memoria}
The map \(\rho_{\rm C}^{\rm disc}\) preserves the identity,
concatenation, and inversion. At complete returns,
\begin{equation}
 \rho_{\rm C}^{\rm disc}(\gamma_{108})=(I,72\ell_0u),\qquad
 \rho_{\rm C}^{\rm disc}(\gamma_{108}^{\,N})=(I,72N\ell_0u).
 \label{vii:eq:traslacion-retorno-108}
\end{equation}
Inversion reverses the sign of the translation.
\end{theorem}

\begin{proof}
The semidirect product is
\((R,a)(S,b)=(RS,a+Rb)\). Since \(R_4u=u\),
\[
\begin{split}
 \rho_{\rm C}^{\rm disc}(\gamma_2)
 \rho_{\rm C}^{\rm disc}(\gamma_1)
 &=
 \bigl(R_4^{c_g(\gamma_2)+c_g(\gamma_1)},
 \ell_0(c_s(\gamma_2)+c_s(\gamma_1))u\bigr)\\
 &=\rho_{\rm C}^{\rm disc}(\gamma_2\gamma_1).
\end{split}
\]
The cocycle of the identity is zero and that of the inverse is its negative.
For the complete circuit, substitute \(c_{108}=(4,72)\) and
\(R_4^4=I\). Concatenation of \(N\) circuits keeps the rotational
part equal to the identity and sums their translations.
\end{proof}

The amplitude \(72\ell_0\) is a length associated with a finite route.
A realization through a cotetrad and connection transports it to its
affine holonomy. The local contribution to the gravitational equations
is then obtained through the curvature of that connection,
variation of the material action, and normalization of its density.
The representation thus preserves the exact cumulative content
that must enter that realization.
